\section{Possibility of Embodiment}

To verify the performance feasibility of the work ''Music for multilayered temporal structures'' described thus far, actual jazz musicians performed ''Etude \#1'' and ''Etude \#2'', which were recorded and documented on video.

The musicians who participated in the performance were three members in a so-called "piano trio" configuration: Tomohito Yoshijima on drums, Yuji Kawamoto on bass, and Masaaki Kaneko on piano. For this performance, the use of existing formats such as "lead sheets" and "piano trio" for the scores and instrumental arrangement was chosen because it was believed to connect to questioning "new music as popular music" rather than being merely a system or experimental approach for this "music with multilayered temporal structures."

The performance was conducted according to the following schedule:

\begin{enumerate}
  \item After determining the musicians, simulation audio tracks and materials were distributed
  \item The first rehearsal was held on 10/22 (Sunday) (drums and bass only)
	\item The second rehearsal was held on 10/29 (Sunday) with piano added
	\item At the second rehearsal, at the request of the bassist and pianist, the Apple Logic Pro X data of the simulation audio was also distributed
	\item Recording and filming were conducted on 11/10 (Friday)
\end{enumerate}

The recording and filming took place in a so-called "recording studio," where the performers were separated into different booths and played while listening to each other's performance through headphones. No click track was used during the recording. This video documentation is publicly available on the author's YouTube account as {\bf Video Documentation ''Music for multilayered temporal structures''}\cite{yamato:04}.

\subsection{Reflections on Rehearsals}

As indicated earlier, two rehearsals were held. The number and participants were determined by the performers' scheduling availability. In the first rehearsal, we worked with the drummer and bassist. Although audio and materials had been distributed in advance, the performers had not specifically practiced for the rehearsal by that day, so we first conducted a reading of the work (booklet) and explained the scores. At that time, the materials did not include explanations of "music with multilayered temporal structures" or "variable meter notation," so explanation by the author was necessary. However, regarding the rehearsal itself, after explaining that "in reality, the author does not exist at this rehearsal" in reference to the work concept, we emphasized the performers' autonomy and endeavored to observe as much as possible. The same approach was taken during the second rehearsal.

In the early stages of rehearsal, both the bassist and drummer tried to proceed with a "fixed meter notation" mindset, but it did not work well. They cited reasons such as "quintuplet quarter notes do not appear in the kind of music they usually play," "when thinking in quintuplet quarter notes, the tuplets cross bar lines when viewed from 4/4 time," and "the downbeat of the quintuplet quarter notes is a rest." Therefore, understanding was facilitated by explaining the structure of the piece as shown in Figure \ref{fig:figure_07} and demonstrating how to feel the rhythm. In this rehearsal, the bassist was unable to leave out the downbeat until the end and started by practicing playing all beats of the quintuplet quarter notes. The drummer proceeded with mastering first the binary division of 5 beats as a reference to understand multiple beats. When moving to ''Etude \#2'' in the latter half of the rehearsal, the transition was smooth, and we received comments from the performers about feeling the effectiveness of ''Etude \#1''.

Next, in the second rehearsal, the pianist joined for a rehearsal toward the actual recording. In this rehearsal, the pianist also initially took a "fixed meter notation" approach, but similarly found it difficult to play quintuplet quarter notes and felt ''Etude \#2'' was more difficult than ''Etude \#1''. Also, as before, after communicating that the author would remain as an observer as much as possible, the rehearsal proceeded, but when performers made structural changes to the music to reduce performance difficulty, this was pointed out and not permitted. For example, since ''Etude \#1'' is made in one mode of D Dorian, the position of rehearsal letter B in the piece was difficult to identify, so referencing {\bf ''So What''}\cite{miles:01}, a request was made to change it up a semitone, but this was not permitted. When actually tried once, while it made the overall structure of the piece easier to understand, we communicated that the purpose of ''Etude \#1'' was "structural understanding of multilayered temporal structures" and should be done in a minimal form. Also, on this day, there was almost no playing finer than quarter notes in the solo sections. This is believed to be because the pianist was trying to understand the quintuplet quarter notes in ''Etude \#2'' with fixed meter notation, converting one beat to five eighth notes in 4/4 time, and thus could not subdivide one beat of 5/4 time.

\subsection{Reflections on Recording and Video Documentation}

Through the two rehearsals, several issues were identified. As mentioned earlier, the bassist and pianist each practiced and approached the recording. ''Etude \#1'' proceeded relatively smoothly, but toward the end of takes, they became unable to complete the performance. Comments indicated this was because the load of each person keeping rhythm was higher than usual pieces. In ''Etude \#2'', it was observed that the pianist independently noticed that "the bass and drum attack points always indicate beats 1 and 3 of the piano." However, regarding the piano solo, the frequency of using notes finer than quarter notes was low, and when pointed out, there were comments that it was difficult to play. The bassist switched rhythms at the B section of ''Etude \#2'' by using the beat of the theme immediately before the B section as a trigger. Particularly in the solo, it can be seen that this part was well utilized.

Throughout the video, it can be observed from the body movements and head movements of the pianist and bassist that each is playing in a different time signature. Also, there were comments such as "during recording (while playing), I don't know what kind of music it is becoming," but on actual playback, "it sounds pop," and "it's like a notated {\bf ''Bitches Brew''}\cite{miles:02}, but I don't know of similar attempts to reproduce this with notation."

Regarding the recording, while there were some innovations and performance compositional ideas (such as the interlude-like development for returning from the solo to the theme in ''Etude \#1'', including the drummer's drum patterns), these did not involve changes to the overall structure of the score and remained at the level of on-site innovations by jazz musicians for performance, while still being able to reproduce the structure of the pieces. Through these rehearsals and recordings, the performance feasibility of "music with multilayered temporal structures" was demonstrated.
